% GENERATED FILE. Edit canonical lessons, then run npm run generate:books.
% canonical-source-hash: fnv1a64:7920a53205b12201
% canonical-lessons: GU-C162-gharkam, GU-C162-ghantadi, GU-C162-sametvu, GU-C162-matlu, GU-C162-khuno

\chapter{Housework and the doorbell}
\label{ch:gu-housework-and-the-doorbell}
\hlchaptermodality{162}

\begin{chapteropening}
\textbf{By the end of this chapter:} \emph{I can talk about housework and things around the house.}

Complete the last lesson of chapter 162: I can talk about housework and things around the house.
\end{chapteropening}

\section[\texorpdfstring{\gu{ઘરકામ} (gharkām)}{ઘરકામ (gharkām)}]{\gu{ઘરકામ} (gharkām) — housework}

\label{lesson:GU-C162-gharkam}



\begin{quote}
Before the new one: say the Gujarati for to break, then the Gujarati for to gather.
\end{quote}

\subsection*{You'll want to know: \gu{ઘરકામ}}
\textbf{\gu{ઘરકામ}} — \emph{gharkām} — \textquotedblleft{}housework\textquotedblright{}.

\begin{quote}
\textbf{\gu{અમે} \gu{સાથે} \gu{ઘરકામ} \gu{કરીએ} \gu{છીએ}.} — \emph{ame sāthe gharkām karīe chīe} — \textquotedblleft{}We do the housework together.\textquotedblright{}
\end{quote}

\begin{cousinweb}
\textbf{\gu{ઘરકામ}} (\emph{gharkām}) is built from two inherited words already taught, \emph{ghar}, \textquotedblleft{}house\textquotedblright{}, and \emph{kām}, \textquotedblleft{}work\textquotedblright{}.
\end{cousinweb}

\begin{grammarlens}[title={What you've built}]
One more word for everyday life at A2.
\end{grammarlens}

\subsection*{Your turn}
\begin{itemize}
\raggedright
  \item \emph{Say it:} \emph{gharkām}
  \item \emph{Say it:} \emph{gharkām}, once more
  \item \emph{Say it:} \emph{ame sāthe gharkām karīe chīe} — We do the housework together.
  \item \emph{Recall:} say the Gujarati for to break, then the Gujarati for to gather, then say \emph{gharkām} again
\end{itemize}

\begin{tcolorbox}[breakable,colback=teal!4,colframe=teal!35!black,title={Before you move on}]
What does \gu{ઘરકામ} mean? (\textquotedblleft{}Housework\textquotedblright{}.) Say it once more.
\end{tcolorbox}
\section[\texorpdfstring{\gu{ઘંટડી} (ghanṭaḍī)}{ઘંટડી (ghanṭaḍī)}]{\gu{ઘંટડી} (ghanṭaḍī) — bell, doorbell}

\label{lesson:GU-C162-ghantadi}



\begin{quote}
Before the new one: say the Gujarati for to gather, then the Gujarati for housework.
\end{quote}

\subsection*{You'll want to know: \gu{ઘંટડી}}
\textbf{\gu{ઘંટડી}} — \emph{ghanṭaḍī} — \textquotedblleft{}bell, doorbell\textquotedblright{}.

\begin{quote}
\textbf{\gu{બારણા} \gu{પાસે} \gu{ઘંટડી} \gu{છે}.} — \emph{bārṇā pāse ghanṭaḍī che} — \textquotedblleft{}There is a bell by the door.\textquotedblright{}
\end{quote}

\begin{cousinweb}
\textbf{\gu{ઘંટડી}} (\emph{ghanṭaḍī}) is a small form built on the word for a bell, which comes from Sanskrit \emph{ghaṇṭā}, \textquotedblleft{}bell\textquotedblright{}. CDIAL lists Old Gujarati \emph{ghāṁṭa}, \textquotedblleft{}bell\textquotedblright{}, but not the Gujarati diminutive itself (CDIAL 4421).
\end{cousinweb}

\begin{grammarlens}[title={What you've built}]
One more word for everyday life at A2.
\end{grammarlens}

\subsection*{Your turn}
\begin{itemize}
\raggedright
  \item \emph{Say it:} \emph{ghanṭaḍī}
  \item \emph{Say it:} \emph{ghanṭaḍī}, once more
  \item \emph{Say it:} \emph{bārṇā pāse ghanṭaḍī che} — There is a bell by the door.
  \item \emph{Recall:} say the Gujarati for to gather, then the Gujarati for housework, then say \emph{ghanṭaḍī} again
\end{itemize}

\begin{tcolorbox}[breakable,colback=teal!4,colframe=teal!35!black,title={Before you move on}]
What does \gu{ઘંટડી} mean? (\textquotedblleft{}Bell, doorbell\textquotedblright{}.) Say it once more.
\end{tcolorbox}
\section[\texorpdfstring{\gu{સમેટવું} (sameṭvũ)}{સમેટવું (sameṭvũ)}]{\gu{સમેટવું} (sameṭvũ) — to tidy up, to fold away}

\label{lesson:GU-C162-sametvu}



\begin{quote}
Before the new one: say the Gujarati for housework, then the Gujarati for bell.
\end{quote}

\subsection*{You'll want to know: \gu{સમેટવું}}
\textbf{\gu{સમેટવું}} — \emph{sameṭvũ} — \textquotedblleft{}to tidy up, to fold away\textquotedblright{}.

\begin{quote}
\textbf{\gu{સવારે} \gu{હું} \gu{પથારી} \gu{સમેટું} \gu{છું}.} — \emph{savāre hũ pathārī sameṭũ chũ} — \textquotedblleft{}In the morning I fold away my bedding.\textquotedblright{}
\end{quote}

\begin{cousinweb}
\textbf{\gu{સમેટવું}} (\emph{sameṭvũ}) is a tadbhava. It comes from Sanskrit \emph{saṁvēṣṭayati}, \textquotedblleft{}wraps up, winds round\textquotedblright{}, and CDIAL glosses the Gujarati verb \textquotedblleft{}to wind up, gather up\textquotedblright{} (CDIAL 13026).
\end{cousinweb}

\begin{grammarlens}[title={What you've built}]
One more word for everyday life at A2.
\end{grammarlens}

\subsection*{Your turn}
\begin{itemize}
\raggedright
  \item \emph{Say it:} \emph{sameṭvũ}
  \item \emph{Say it:} \emph{sameṭvũ}, once more
  \item \emph{Say it:} \emph{savāre hũ pathārī sameṭũ chũ} — In the morning I fold away my bedding.
  \item \emph{Recall:} say the Gujarati for housework, then the Gujarati for bell, then say \emph{sameṭvũ} again
\end{itemize}

\begin{tcolorbox}[breakable,colback=teal!4,colframe=teal!35!black,title={Before you move on}]
What does \gu{સમેટવું} mean? (\textquotedblleft{}To tidy up, to fold away\textquotedblright{}.) Say it once more.
\end{tcolorbox}
\section[\texorpdfstring{\gu{માટલું} (māṭlũ)}{માટલું (māṭlũ)}]{\gu{માટલું} (māṭlũ) — clay pot, water pot}

\label{lesson:GU-C162-matlu}



\begin{quote}
Before the new one: say the Gujarati for bell, then the Gujarati for to tidy up.
\end{quote}

\subsection*{You'll want to know: \gu{માટલું}}
\textbf{\gu{માટલું}} — \emph{māṭlũ} — \textquotedblleft{}clay pot, water pot\textquotedblright{}.

\begin{quote}
\textbf{\gu{માટલામાં} \gu{પાણી} \gu{છે}.} — \emph{māṭlāmā̃ pāṇī che} — \textquotedblleft{}There is water in the clay pot.\textquotedblright{}
\end{quote}

\begin{cousinweb}
\textbf{\gu{માટલું}} (\emph{māṭlũ}) is a tadbhava. It comes from \emph{*mārtta}, \textquotedblleft{}earthen\textquotedblright{}, beside Sanskrit \emph{mārttika}, \textquotedblleft{}earthen pot\textquotedblright{}, and CDIAL glosses the Gujarati word \textquotedblleft{}large earthen pot\textquotedblright{} (CDIAL 10085).
\end{cousinweb}

\begin{grammarlens}[title={What you've built}]
One more word for everyday life at A2.
\end{grammarlens}

\subsection*{Your turn}
\begin{itemize}
\raggedright
  \item \emph{Say it:} \emph{māṭlũ}
  \item \emph{Say it:} \emph{māṭlũ}, once more
  \item \emph{Say it:} \emph{māṭlāmā̃ pāṇī che} — There is water in the clay pot.
  \item \emph{Recall:} say the Gujarati for bell, then the Gujarati for to tidy up, then say \emph{māṭlũ} again
\end{itemize}

\begin{tcolorbox}[breakable,colback=teal!4,colframe=teal!35!black,title={Before you move on}]
What does \gu{માટલું} mean? (\textquotedblleft{}Clay pot, water pot\textquotedblright{}.) Say it once more.
\end{tcolorbox}
\section[\texorpdfstring{\gu{ખૂણો} (khūṇo)}{ખૂણો (khūṇo)}]{\gu{ખૂણો} (khūṇo) — corner}

\label{lesson:GU-C162-khuno}



\begin{quote}
Before the new one: say the Gujarati for to tidy up, then the Gujarati for clay pot.
\end{quote}

\subsection*{You'll want to know: \gu{ખૂણો}}
\textbf{\gu{ખૂણો}} — \emph{khūṇo} — \textquotedblleft{}corner\textquotedblright{}.

\begin{quote}
\textbf{\gu{માટલું} \gu{રસોડાના} \gu{ખૂણામાં} \gu{છે}.} — \emph{māṭlũ rasoḍānā khūṇāmā̃ che} — \textquotedblleft{}The clay pot is in the corner of the kitchen.\textquotedblright{}
\end{quote}

\begin{cousinweb}
\textbf{\gu{ખૂણો}} (\emph{khūṇo}) is a tadbhava. It comes from \emph{*khuṇṭa}, \textquotedblleft{}corner\textquotedblright{}, a reconstructed form, crossed with Sanskrit \emph{koṇa}, \textquotedblleft{}corner\textquotedblright{}; CDIAL lists the Gujarati \emph{khūṇo}, \textquotedblleft{}corner\textquotedblright{} (CDIAL 3898).
\end{cousinweb}

\begin{grammarlens}[title={What you've built}]
One more word for everyday life at A2.
\end{grammarlens}

\subsection*{Your turn}
\begin{itemize}
\raggedright
  \item \emph{Say it:} \emph{khūṇo}
  \item \emph{Say it:} \emph{khūṇo}, once more
  \item \emph{Say it:} \emph{māṭlũ rasoḍānā khūṇāmā̃ che} — The clay pot is in the corner of the kitchen.
  \item \emph{Recall:} say the Gujarati for to tidy up, then the Gujarati for clay pot, then say \emph{khūṇo} again
\end{itemize}

\begin{tcolorbox}[breakable,colback=teal!4,colframe=teal!35!black,title={Before you move on}]
What does \gu{ખૂણો} mean? (\textquotedblleft{}Corner\textquotedblright{}.) Say it once more.
\end{tcolorbox}
